\section{Доказательство теоремы Галуа о разрешимости}

Напомним определение (см.\ перед \rf{3.28}). Многочлен $f\in F[X]$ (или уравнение $f(X)=0$) \emph{разрешим в радикалах},
если существует башня полей
\[
F=F_0\subset F_1\subset F_2\subset\dots\subset F_m
\]
такая, что (a) $F_i=F_{i-1}[\alpha_i]$, где $\alpha_i^{r_i}\in F_{i-1}$ для некоторого целого $r_i\ge1$, и (b) $F_m$ содержит
поле разложения $f$. Иначе говоря, корни $f$ получаются из элементов $F$ сложением, вычитанием, умножением, делением и
извлечением корней. Такую башню называем \emph{радикальной}. Группа $G$ \emph{разрешима}, если у неё есть цепочка
подгрупп $G=G_0\supset G_1\supset\dots\supset G_m=\{1\}$, где каждая $G_i$ нормальна в $G_{i-1}$ и $G_{i-1}/G_i$ абелева;
для конечной разрешимой группы цепочку можно уплотнить так, что все факторы $G_{i-1}/G_i$ будут циклическими (даже
простого порядка). Подгруппы и факторгруппы разрешимых групп разрешимы.

Для многочлена $f$ над полем характеристики $0$ его поле разложения $E_f$~--- это поле разложения произведения различных
неприводимых множителей $f$, а этот многочлен сепарабелен; поэтому $E_f/F$~--- расширение Галуа, и группа $G_f=\Gal(E_f/F)$
определена для любого $f$.

\begin{lemma}[5.33]
Пусть $f\in F[X]$ сепарабелен и $F'$~--- поле, содержащее $F$. Тогда группа Галуа $f$ как многочлена из $F'[X]$
изоморфна подгруппе группы Галуа $f$ как многочлена из $F[X]$.
\end{lemma}

\begin{proof}
Пусть $E'$~--- поле разложения $f$ над $F'$ и $\alpha_1,\dots,\alpha_m$~--- корни $f$ в $E'$. Тогда $E=F[\alpha_1,\dots,\alpha_m]$~---
поле разложения $f$ над $F$ (оно порождено корнями, и в нём $f$ раскладывается на линейные множители). Каждый
$\sigma\in\Gal(E'/F')$ переставляет корни $\alpha_i$ (он сохраняет коэффициенты $f$) и потому переводит $E$ в себя; он
тождествен на $F'\supset F$, так что $\sigma|_E\in\Gal(E/F)$. Отображение $\sigma\mapsto\sigma|_E$~--- гомоморфизм групп.
Он инъективен: если $\sigma|_E=\id$, то $\sigma$ оставляет на месте все $\alpha_i$, а $E'=F'[\alpha_1,\dots,\alpha_m]$, так
что $\sigma=\id$.
\end{proof}

\[
\begin{tikzcd}[column sep=large,row sep=large]
F'\arrow[r,hook] & E'=F'[\alpha_1,\dots,\alpha_m]\\
F\arrow[u,hook]\arrow[r,hook] & E=F[\alpha_1,\dots,\alpha_m]\arrow[u,hook,cnew]
\end{tikzcd}
\qquad\qquad
\begin{tikzcd}
\Gal(E'/F')\arrow[r,hook,cnew,"\sigma\,\mapsto\,\sigma|_E"] & \Gal(E/F)
\end{tikzcd}
\]
\figcap{Лемма \rf{5.33}: поле разложения $E$ над меньшим полем лежит в поле разложения $E'$ над большим; ограничение на $E$
вкладывает $\Gal(E'/F')$ в $\Gal(E/F)$. Например, при $F'\supset E$ группа $\Gal(E'/F')$ тривиальна.}

\begin{theorem}[5.34]
Пусть $F$~--- поле характеристики $0$. Многочлен $f\in F[X]$ разрешим в радикалах тогда и только тогда, когда его группа
Галуа $G_f$ разрешима.
\end{theorem}

Идея доказательства. Разрешимость группы означает, что расширение можно разбить на шаги с циклическими группами, а по
теореме \rf{5.27} циклический шаг~--- это присоединение радикала, если в поле есть нужные корни из единицы. Поэтому сначала
присоединяют корни из единицы. Обратно, после присоединения корней из единицы радикальная башня становится цепочкой
абелевых шагов, но вся башня может не быть расширением Галуа; поэтому её достраивают до нормального замыкания,
порождённого всеми сопряжёнными радикалами.

\begin{proof}
Заменив $f$ произведением его различных неприводимых множителей (поле разложения и группа Галуа от этого не
меняются), можно считать $f$ сепарабельным.

$\Leftarrow$: Пусть $G_f$ разрешима, $M=(\deg f)!$, и $\zeta$~--- первообразный корень степени $M$ из единицы (в
некотором расширении $F$; он существует по \rf{5.8}, так как $\operatorname{char}F=0$). Положим $F'=F[\zeta]$, и пусть $E$~---
поле разложения $f$ над $F'$, $G=\Gal(E/F')$. По лемме \rf{5.33} $G$ изоморфна подгруппе $G_f$ и потому разрешима. Возьмём
цепочку
\[
G=G_0\supset G_1\supset\dots\supset G_m=\{1\},\qquad G_i\ \text{нормальна в}\ G_{i-1},\quad G_{i-1}/G_i\ \text{циклична},
\]
и положим $F_i=E^{G_i}$. Получаем башню
\[
F\subset F[\zeta]=F'=F_0\subset F_1\subset\dots\subset F_m=E.
\]
По основной теореме (\rf{3.17}) $E/F_{i-1}$~--- расширение Галуа с группой $G_{i-1}$, а так как $G_i$ нормальна в $G_{i-1}$,
расширение $F_i/F_{i-1}$~--- расширение Галуа с группой $G_{i-1}/G_i$, т.\,е.\ циклическое некоторой степени $d_i$. Число $d_i$
делит $|G|$, а $|G|$ делит $|G_f|$, которое делит $(\deg f)!=M$ ($G_f$ вкладывается в группу перестановок корней). Поэтому
$\zeta^{M/d_i}\in F_{i-1}$~--- первообразный корень степени $d_i$ из единицы, и по \rf{5.27} (над полем $F_{i-1}$)
$F_i=F_{i-1}[\alpha_i]$ с $\alpha_i^{d_i}\in F_{i-1}$. Первый шаг $F\subset F[\zeta]$ тоже радикальный: $\zeta^M=1\in F$. Наконец,
$E$ содержит поле разложения $f$ над $F$ (подполе, порождённое над $F$ корнями $f$). Итак, $f$ разрешим в радикалах.

\begin{center}
\begin{tikzpicture}[every node/.style={font=\small}]
\node (G0) at (0,0) {$G=G_0$}; \node (G1) at (0,1.2) {$G_1$}; \node (Gd) at (0,2.4) {$\vdots$}; \node (Gm) at (0,3.6) {$G_m=\{1\}$};
\node at (0,0.6) {\rotatebox{-90}{$\subset$}}; \node at (0,1.8) {\rotatebox{-90}{$\subset$}}; \node at (0,3.0) {\rotatebox{-90}{$\subset$}};
\node[cold] at (-1.6,0.6) {$G_0/G_1$ циклична};
\node (F0) at (6,0) {$F'=F[\zeta]=F_0$}; \node (F1) at (6,1.2) {$F_1=F_0[\alpha_1]$}; \node at (6,2.4) {$\vdots$};
\node (Fm) at (6,3.6) {$F_m=E$}; \node (F) at (6,-1.2) {$F$};
\node at (6,0.6) {\rotatebox{90}{$\subset$}}; \node at (6,1.8) {\rotatebox{90}{$\subset$}}; \node at (6,3.0) {\rotatebox{90}{$\subset$}};
\node at (6,-0.6) {\rotatebox{90}{$\subset$}};
\node[cnew] at (8.4,0.6) {$\alpha_1^{d_1}\in F_0$ (5.27)}; \node[cnew] at (8.1,-0.6) {$\zeta^{M}=1$};
\draw[<->,gray] (1.2,0) -- (4.6,0); \draw[<->,gray] (0.6,1.2) -- (4.9,1.2); \draw[<->,gray] (1.2,3.6) -- (5.2,3.6);
\node[gray] at (3,3.85) {$F_i=E^{G_i}$};
\end{tikzpicture}
\end{center}
\figcap{Доказательство <<$\Leftarrow$>>: цепочке подгрупп с циклическими факторами отвечает башня циклических расширений;
так как в $F'$ уже есть все нужные корни из единицы, каждый шаг~--- присоединение радикала.}

$\Rightarrow$: Достаточно показать, что $G_f$~--- факторгруппа разрешимой группы, т.\,е.\ найти расширение Галуа
$\tilde E\supset F$ с разрешимой группой, содержащее поле разложения $E$ многочлена $f$: тогда $G_f=\Gal(E/F)\cong
\Gal(\tilde E/F)/\Gal(\tilde E/E)$ (\rf{3.17}) разрешима. На языке групп: нужно представить $G_f$ как факторгруппу
разрешимой группы по нормальной подгруппе; на языке полей: $F\subset E\subset\tilde E$, все расширения~--- Галуа.

Пусть дана радикальная башня $F=F_0\subset\dots\subset F_m$, $F_i=F_{i-1}[\alpha_i]$, $\alpha_i^{r_i}\in F_{i-1}$, причём
$F_m\supset E$. Положим $n=r_1\cdots r_m$. Выберем расширение Галуа $\Om/F$, содержащее (копию) $F_m$ и первообразный
корень $\zeta$ степени $n$ из единицы. Например: расширение $F_m/F$ сепарабельно (характеристика $0$), поэтому
$F_m=F[\gamma]$ (\rf{5.1}); пусть $g$~--- минимальный многочлен $\gamma$ над $F$, и $\Om$~--- поле разложения $g(X)(X^n-1)$ над
$F$. Оно является расширением Галуа (характеристика $0$), содержит корень $\gamma'$ многочлена $g$ и $F$-изоморфное $F_m$
поле $F[\gamma']$, а также первообразный корень степени $n$ из единицы (\rf{5.8}). Перенеся башню в $\Om$ этим изоморфизмом,
можно считать, что $F_m\subset\Om$.

Расширение $F_m/F$ может не быть расширением Галуа (например, $\Q\subset\Q[\sqrt[3]2]$). Поэтому перейдём к нормальному
замыканию. Пусть $G=\Gal(\Om/F)$ и $H=\Gal(\Om/F_m[\zeta])$, так что $F_m[\zeta]=\Om^H$. Наибольшая нормальная подгруппа
$G$, содержащаяся в $H$,~--- это $N=\bigcap_{\sigma\in G}\sigma H\sigma^{-1}$, и по \rf{3.18}(b) поле
\[
\tilde E=\Om^N=\prod_{\sigma\in G}\sigma\bigl(F_m[\zeta]\bigr)\qquad\text{(композит)}
\]
~--- наименьшее расширение Галуа поля $F$ в $\Om$, содержащее $F_m[\zeta]$ (нормальное замыкание).

\begin{center}
\begin{tikzpicture}[every node/.style={font=\small},yscale=0.95]
\node (F) at (0,0) {$F$};
\node (E) at (-1.6,1.3) {$E$};
\node (Fm) at (0,2.5) {$F_m$};
\node (Fmz) at (0,3.6) {$F_m[\zeta]=\Om^H$};
\node (Et) at (0,4.9) {$\tilde E=\Om^N$};
\node (Om) at (0,6.2) {$\Om$};
\draw (F) -- (E) -- (Fm); \draw (F) -- node[right,cold,align=left]{радикальная\\ башня} (Fm); \draw (Fm) -- (Fmz); \draw (Fmz) -- (Et); \draw (Et) -- (Om);
\node[cold] at (-3.4,1.3) {поле разложения $f$};
\node (G) at (6,0) {$G=\Gal(\Om/F)$};
\node (H) at (6,3.6) {$H$};
\node (N) at (6,4.9) {$N=\bigcap_{\sigma\in G}\sigma H\sigma^{-1}$};
\node (one) at (6,6.2) {$\{1\}$};
\draw (G) -- (H) -- (N) -- (one);
\draw[<->,gray,dashed] (Fmz.east) -- (H.west);
\draw[<->,gray,dashed] (Et.east) -- (N.west);
\draw[decorate,decoration={brace,mirror,amplitude=5pt},cnew] (7.75,0) -- node[right=6pt,align=left]{$N$ нормальна в $G$,\\
 $G/N\cong\Gal(\tilde E/F)$\\ разрешима} (7.75,4.9);
\end{tikzpicture}
\end{center}
\figcap{Расширение $F_m/F$ может не быть расширением Галуа, поэтому его заменяют нормальным замыканием
$\tilde E\supset F_m[\zeta]$; ему отвечает наибольшая нормальная подгруппа $N\subset H$. Поле разложения $E$ лежит в
$F_m\subset\tilde E$, и $G_f$~--- факторгруппа $\Gal(\tilde E/F)$.}

Так как $F_m[\zeta]=F[\zeta,\alpha_1,\dots,\alpha_m]$ и $\sigma\zeta$~--- тоже первообразный корень степени $n$, т.\,е.\ степень
$\zeta$, поле $\tilde E$ порождено над $F$ элементами
\[
\zeta;\quad\alpha_1,\dots,\alpha_m;\quad\sigma_2\alpha_1,\dots,\sigma_2\alpha_m;\quad\dots;\quad\sigma_s\alpha_1,\dots,\sigma_s\alpha_m,
\]
где $\sigma_1=1,\sigma_2,\dots,\sigma_s$~--- все элементы $G$. Присоединим эти элементы к $F$ по одному в указанном
порядке:
\[
F\subset F[\zeta]\subset F[\zeta,\alpha_1]\subset\dots\subset F_m[\zeta]\subset F_m[\zeta,\sigma_2\alpha_1]\subset\dots\subset\tilde E.
\]
Каждый шаг $K\subset K[\beta]$ этой башни~--- присоединение корня: при присоединении $\beta=\sigma_j\alpha_i$ имеем
\[
\beta^{r_i}=\sigma_j\bigl(\alpha_i^{r_i}\bigr)\in\sigma_j(F_{i-1})=F[\sigma_j\alpha_1,\dots,\sigma_j\alpha_{i-1}]\subset K,
\]
потому что $\sigma_j\alpha_1,\dots,\sigma_j\alpha_{i-1}$ присоединены раньше (именно для этого важен порядок). Все шаги~---
абелевы расширения Галуа:
\begin{itemize}
\item $F[\zeta]/F$~--- поле разложения $X^n-1$, его группа вкладывается в $(\Z/n\Z)^\times$ (\rf{5.8});
\item если $K\supset F[\zeta]$ и $\beta^r=b\in K$, $r\mid n$, то при $b=0$ шаг тривиален, а при $b\ne0$ корни $X^r-b$~--- это
$r$ различных элементов $\zeta^{jn/r}\beta\in K[\beta]$; значит, $K[\beta]$~--- поле разложения сепарабельного многочлена, и
$\tau\mapsto\tau\beta/\beta$~--- инъективный гомоморфизм $\Gal(K[\beta]/K)$ в циклическую группу $\mu_r$ (как перед \rf{5.27}).
\end{itemize}
Если $F=K_0\subset K_1\subset\dots\subset K_t=\tilde E$~--- эта башня, то подгруппы $\Gal(\tilde E/K_j)$ образуют цепочку,
$\Gal(\tilde E/K_j)$ нормальна в $\Gal(\tilde E/K_{j-1})$ (так как $K_j/K_{j-1}$~--- расширение Галуа), и факторгруппа
$\Gal(\tilde E/K_{j-1})/\Gal(\tilde E/K_j)\cong\Gal(K_j/K_{j-1})$ абелева (\rf{3.17}). Значит, $\Gal(\tilde E/F)$ разрешима. Так как
$E\subset F_m\subset\tilde E$ и $E/F$~--- расширение Галуа, группа $G_f=\Gal(E/F)$~--- факторгруппа разрешимой группы
$\Gal(\tilde E/F)$ и потому разрешима.
\end{proof}

В доказательстве характеристика $0$ использовалась дважды: для сепарабельности (существование $\Om$ и $G_f$) и для
существования первообразных корней из единицы любой степени.

\begin{aside}[5.35]
Одно из главных достижений Галуа~--- следующая теорема: неприводимый многочлен простой степени из $\Q[X]$ разрешим в
радикалах тогда и только тогда, когда его поле разложения порождается любыми двумя его корнями. Например, у $X^5-2$
корни~--- это $\zeta^i\alpha$, $\alpha=\sqrt[5]2$, $\zeta=\zeta_5$; поле, порождённое двумя корнями $\zeta^i\alpha$ и $\zeta^j\alpha$
($i\ne j$), содержит $\zeta^{j-i}$, а значит, и $\zeta$, и $\alpha$, т.\,е.\ совпадает с полем разложения $\Q[\zeta,\alpha]$ степени
$20$. А у многочлена $X^5-6X+3$ с группой $S_5$ (\rf{5.35e}) два корня
порождают поле степени $120/3!=20$ (неподвижное поле стабилизатора двух корней), меньшее поля разложения степени $120$.
\end{aside}

\subsection*{Разрешимость в характеристике $p$}

В характеристике $p$ теорема \rf{5.34} в прежней формулировке неверна по двум причинам (ср.\ \rf{3.29}): многочлен может быть
несепарабельным, и, главное, циклическое расширение степени $p$ не обязано получаться присоединением радикалов: по
\rf{5.32j} уравнение $X^p-X-t=0$ над $\F_p(t)$ имеет группу Галуа $\Z/p\Z$, но не разрешимо в радикалах. Исправление
подсказывает теория Артина--Шрайера: к радикалам нужно добавить корни многочленов $X^p-X-a$.

\begin{definition}[5.35a]
Пусть $\operatorname{char}F=p>0$. Башня $F=F_0\subset F_1\subset\dots\subset F_m$ называется \emph{обобщённой радикальной},
если $F_i=F_{i-1}[\alpha_i]$, где для каждого $i$ либо $\alpha_i^{r_i}\in F_{i-1}$ для некоторого $r_i\ge1$, не делящегося на
$p$, либо $\alpha_i^p-\alpha_i\in F_{i-1}$. Многочлен $f\in F[X]$ \emph{разрешим в обобщённых радикалах}, если $F_m$ содержит
поле разложения $f$ для некоторой обобщённой радикальной башни.
\end{definition}

Корни $p$-й степени мы не разрешаем: они дают чисто несепарабельные шаги. Все шаги обобщённой радикальной башни
сепарабельны: при $\alpha_i^{r_i}=b\ne0$ многочлен $X^{r_i}-b$ сепарабелен (его производная $r_iX^{r_i-1}$ имеет лишь корень
$0$, так как $p\nmid r_i$), при $b=0$ шаг тривиален, а $X^p-X-b$ сепарабелен по \rf{5.32c}(a).

\begin{theorem}[5.35b: теорема Галуа в характеристике $p$]
Пусть $\operatorname{char}F=p>0$ и $f\in F[X]$ сепарабелен. Тогда $f$ разрешим в обобщённых радикалах тогда и только
тогда, когда группа $G_f$ разрешима.
\end{theorem}

\begin{proof}
Доказательство повторяет доказательство \rf{5.34}; новые места~--- шаги степени $p$ и отсутствие корней из единицы степени,
делящейся на $p$.

$\Leftarrow$: Пусть $|G_f|=p^em$, $p\nmid m$, $\zeta$~--- первообразный корень степени $m$ из единицы (\rf{5.8}: $p\nmid m$),
$F'=F[\zeta]$, $E$~--- поле разложения $f$ над $F'$, $G=\Gal(E/F')$. По \rf{5.33} $G$ разрешима; возьмём цепочку
$G=G_0\supset\dots\supset G_k=\{1\}$, $G_i$ нормальна в $G_{i-1}$, с факторами $G_{i-1}/G_i$ \emph{простого} порядка $\ell_i$, и
$F_i=E^{G_i}$. Как в \rf{5.34}, $F_i/F_{i-1}$~--- циклическое расширение степени $\ell_i$, и $\ell_i$ делит $|G_f|$.
\begin{itemize}
\item Если $\ell_i\ne p$, то $\ell_i\mid m$, $\zeta^{m/\ell_i}\in F_{i-1}$~--- первообразный корень степени $\ell_i$ из единицы, и
по \rf{5.27} $F_i=F_{i-1}[\alpha_i]$, $\alpha_i^{\ell_i}\in F_{i-1}$, $p\nmid\ell_i$.
\item Если $\ell_i=p$, то по \rf{5.32c}(c) $F_i=F_{i-1}[\alpha_i]$, $\alpha_i^p-\alpha_i\in F_{i-1}$.
\end{itemize}
Первый шаг $F\subset F[\zeta]$ допустим: $\zeta^m=1\in F$, $p\nmid m$. Значит, $F\subset F'=F_0\subset\dots\subset F_k=E$~---
обобщённая радикальная башня, и $E$ содержит поле разложения $f$ над $F$.

$\Rightarrow$: Пусть $F=F_0\subset\dots\subset F_m$~--- обобщённая радикальная башня, $F_m\supset E$ (поле разложения $f$
над $F$), и $n$~--- произведение тех $r_i$, которые встречаются в шагах первого типа ($n=1$, если таких нет); $p\nmid n$.
Каждый $\alpha_i$ сепарабелен над $F_{i-1}$, т.\,е.\ $F_i/F_{i-1}$ сепарабельно (\rf{3.15}), поэтому по транзитивности
сепарабельности (\rf{6.8b}) расширение $F_m/F$ сепарабельно, и дальше всё идёт, как в \rf{5.34}: $F_m=F[\gamma]$ (\rf{5.1}), $\Om$~--- поле разложения над $F$ произведения различных
неприводимых множителей многочлена $g(X)(X^n-1)$ ($g$~--- минимальный многочлен $\gamma$; все эти множители
сепарабельны, так как $p\nmid n$), так что $\Om/F$~--- расширение Галуа, содержащее копию $F_m$ и первообразный корень
$\zeta$ степени $n$ из единицы. Нормальное замыкание $\tilde E$ поля $F_m[\zeta]$ в $\Om$ порождено элементами $\zeta$ и
$\sigma\alpha_i$, $\sigma\in G=\Gal(\Om/F)$, и, присоединяя их в том же порядке, что в \rf{5.34}, получаем башню
$F=K_0\subset\dots\subset K_t=\tilde E$, где $\sigma_j\alpha_i$ удовлетворяет над предыдущим полем $K$ уравнению
$X^{r_i}=\sigma_j(\alpha_i^{r_i})$ или $X^p-X=\sigma_j(\alpha_i^p-\alpha_i)$ с правой частью из $\sigma_j(F_{i-1})\subset K$. Шаги
первого типа и шаг $F[\zeta]/F$~--- абелевы расширения Галуа, как в \rf{5.34}, а шаги второго типа по \rf{5.32c}(b) либо
тривиальны, либо циклические степени $p$. Значит, $\Gal(\tilde E/F)$ разрешима, и $G_f$~--- её факторгруппа.
\end{proof}

Условие сепарабельности $f$ в \rf{5.35b} существенно: все шаги обобщённой радикальной башни сепарабельны, поэтому (\rf{6.8b})
поле $F_m$ сепарабельно над $F$ и не содержит корней несепарабельных неприводимых многочленов. Например, многочлен
$X^p-t$ над $\F_p(t)$ (он неприводим по \rf{2.3} и несепарабелен, так как его производная равна нулю) в обобщённых радикалах
не разрешим, хотя в обычных разрешим ($\alpha^p=t$). А $X^p-X-t$, наоборот, разрешим в обобщённых радикалах, но не в
обычных (\rf{5.32j}). Так что эти два понятия несравнимы.

\begin{remark}[5.35c]
(a) Над конечным полем $\F_q$ каждый многочлен разрешим даже в обычных радикалах: поле разложения $\F_{q^r}$ порождено над
$\F_q$ образующей циклической группы $\F_{q^r}^\times$, т.\,е.\ корнем из единицы степени $q^r-1$ (взаимно простой с $p$).
Например, корни $X^p-X-1$ над $\F_p$~--- корни из единицы степени $\frac{p^p-1}{p-1}$ (\rf{5.32g}). Так что пример \rf{5.32j}
требует бесконечного поля.

Это не противоречит \rf{5.32i}: там основное поле должно содержать первообразные корни степени $\ell$ из единицы для всех
простых $\ell\ne p$, делящих показатели, а здесь это не так. Пусть $\F_{q^r}$ содержит корни неприводимого многочлена
Артина--Шрайера, т.\,е.\ $\F_{q^p}\subset\F_{q^r}$ и $p\mid r$. Тогда показатель $q^r-1$ делится на
$Q=\frac{q^p-1}{q-1}=1+q+\dots+q^{p-1}>1$, а ни один простой делитель $\ell$ числа $Q$ не делит $q-1$: при
$q\equiv1\pmod\ell$ было бы $Q\equiv p\pmod\ell$, т.\,е.\ $\ell=p$, но $Q\equiv1\pmod p$. Значит, в $\F_q$ нет
первообразных корней степени $\ell$ из единицы. Например, при $q=p=5$ корни $X^5-X-1$~--- корни из единицы степени
$781=11\cdot71$, а в $\F_5$ нет корней степени $11$ и $71$ из единицы, кроме $1$.

(b) Пусть $F$ содержит первообразные корни из единицы всех простых степеней $\ell\ne p$ (например, $F=\Fpb(t)$). По \rf{5.32i}
степень любого расширения Галуа $L/F$, содержащегося в радикальной башне над $F$, не делится на $p$. Поэтому
сепарабельный многочлен $f$, у которого $p$ делит $|G_f|=[E_f:F]$, не разрешим над $F$ в радикалах (даже если $G_f$
разрешима и $f$ разрешим в обобщённых радикалах). Наоборот, если $G_f$ разрешима и $p\nmid|G_f|$, то в доказательстве \rf{5.35b} шагов Артина--Шрайера не возникает, и
$f$ разрешим в обычных радикалах.
\end{remark}

\subsection*{Уравнения степеней 3, 4 и 5}

\begin{aside}[5.35d]
Пусть $\operatorname{char}F=0$ и $f\in F[X]$ степени $n\le4$. Группа $G_f$ вкладывается в $S_n$, а группы $S_2$, $S_3$ и $S_4$
разрешимы: $S_3\supset A_3\supset\{1\}$ с факторами $\Z/2$ и $\Z/3$, $S_4\supset A_4\supset V\supset\{1\}$ ($V$~--- четверная
группа Клейна) с факторами $\Z/2$, $\Z/3$, $(\Z/2)^2$. По \rf{5.34} \emph{каждое уравнение степени не выше $4$ разрешимо в
радикалах}~--- это формулы Кардано и Феррари. Доказательство \rf{5.34} подсказывает и сами формулы: для кубического
многочлена шаг $A_3\subset S_3$ отвечает присоединению квадратного корня из дискриминанта (\rf{4.7}), после чего остаётся
циклическое расширение степени $3$, т.\,е.\ (после присоединения $\zeta_3$) извлечение кубического корня из резольвенты
Лагранжа; для многочлена степени $4$ подгруппе $V$ отвечает поле разложения резольвентной кубики (\rf{4.8}, \rf{4.9}), и задача
сводится к кубической и затем к квадратичным.

При $n=5$ группа $S_5$ неразрешима: её единственная собственная нетривиальная нормальная подгруппа $A_5$ проста и
неабелева. Поэтому многочлен с группой $S_5$ не разрешим в радикалах; такие многочлены есть, например, по \rf{4.16}, \rf{4.17}.
Многочлены пятой степени с разрешимой группой разрешимы в радикалах.
\end{aside}

\begin{example}[5.35e]
\emph{$X^5-6X+3$ над $\Q$ не разрешим в радикалах.} Он неприводим по признаку Эйзенштейна (\rf{1.16}) для $p=3$. Значения в
точках $-2,-1,1,2$ равны $-17,\,8,\,-2,\,23$, поэтому у него не меньше трёх вещественных корней; а производная $5X^4-6$ имеет
лишь два вещественных корня, так что по теореме Ролля вещественных корней не больше трёх. Значит, ровно два корня
невещественны, и по \rf{4.16} $G_f=S_5$, т.\,е.\ по \rf{5.34} уравнение $X^5-6X+3=0$ не разрешимо в радикалах. (Так же $G_f=S_5$ у
$X^5-X-1$, \rf{4.30}, и у $X^5-5X^3+4X-1$ с пятью вещественными корнями, \rf{4.18}.)

\emph{$X^5-2$ разрешим в радикалах}~--- его корни $\zeta_5^k\sqrt[5]2$ записаны радикалами. Согласованно с \rf{5.34} его группа
Галуа (\rf{3.23}) имеет порядок $20$ и разрешима: $\Gal(\Q[\zeta_5,\sqrt[5]2]/\Q[\zeta_5])\cong\Z/5$~--- нормальная подгруппа с
факторгруппой $\Gal(\Q[\zeta_5]/\Q)\cong(\Z/5\Z)^\times\cong\Z/4$.
\end{example}
